\documentclass[11pt]{article}

\usepackage[margin=1in]{geometry}
\usepackage{amsmath,amssymb,amsthm,mathtools}
\usepackage{enumitem}
\usepackage{booktabs}
\usepackage{tabularx}
\usepackage{array}
\usepackage{microtype}
\usepackage[hidelinks]{hyperref}
\usepackage[nameinlink,noabbrev]{cleveref}
\usepackage{xurl}

\newtheorem{definition}{Definition}
\newtheorem{proposition}{Proposition}
\newtheorem{remark}{Remark}
\newcolumntype{Y}{>{\raggedright\arraybackslash}X}

\newcommand{\eps}{\varepsilon}
\newcommand{\defeq}{\stackrel{\mathrm{def}}{=}}

\title{Receipt Accounting and Horizon Composition for Anonymous-DHT Budgets:\\
\large A Rulebook for Repeats, Retries, and State (Series Synthesis)}
\author{}
\date{Unified archive note v0.06 (2026-03-23; archive v651)}

\begin{document}
\maketitle

\begin{abstract}
This archive publishes anonymity budgets as small, replayable receipt line items.
A recurring failure mode in the literature and in deployments is \emph{silent horizon drift}: a proof that is per-lookup is later quoted as ``per day'' or ``for $Q$ lookups'', even though retries, caches, and adaptive routing change what composes.
This note is the series' compact rulebook: it states what a receipt must declare before it can honestly claim horizon-$n$ composition, and it points to the owning papers for proofs.
The goal is editorial: keep mechanism papers short by making ``how budgets compose across time and state'' a single citable module.
\end{abstract}

\paragraph{Scope.}
This note concerns \emph{repetition}: how per-lookup budgets become horizon budgets under retries, windowing, and state evolution.
It does not introduce new endpoint metrics; it only standardizes when a receipt may claim ``$n$-use'' anonymity.

\paragraph{Imports.}
Sequential/adaptive composition for posterior-mass inflation / PML and $(\eps,\delta)$ checkpoints is owned by Anonymity~A.\cite{series:oddsinflation}
Threat-window declarations (unit, window, reset semantics) are owned by Synthesis~4.\cite{series:threatwindows}
Receipt line-item schemas (including plan replay hooks) are owned by Synthesis~12.\cite{series:receiptitems}
Probabilistic separation adjuncts for primary-path guards (a $(0,\rho)$ checkpoint) are owned by Synthesis~23.\cite{series:probsep}
State evolution as a witness surface is owned by the AnonDHT State series.\cite{series:state1}

\paragraph{Later-terse-paper citation posture.}
When a later short paper only needs to answer \emph{what exact public packet makes this horizon-$Q$ claim honest?}, it should stop at the minimal public horizon-claim card below, the compact first-reopen-owner card, and the tiny composition drill.
It should widen only when the live issue becomes the sequential-composition theorem root (Anonymity~A), threat-window semantics (Synthesis~4), the receipt tuple and replay hooks (Synthesis~12), conditional per-step certificates (Synthesis~21), plan/state equivalence (Synthesis~18), or the mechanism-specific state / retry owners that justify a particular regime choice.\cite{series:oddsinflation,series:threatwindows,series:receiptitems,series:condcert,series:planstateequiv,series:operationalA,series:operationalB,series:state1}

\begin{table}[t]
\centering
\small
\begin{tabularx}{0.97\linewidth}{@{}p{0.30\linewidth}Y@{}}
\toprule
\textbf{Public row} & \textbf{Pinned value}\\
\midrule
Threat window & \texttt{tw\_id = lookup.epoch.v1}; one lookup epoch is the unit, with reset only at the declared lookup boundary \\
Regime declaration & Regime~II adaptive sequential composition over $Q=20$ uses; the receipt therefore does not rely on independence and must publish a replayable conditional-step packet \\
Per-unit budget & \texttt{budget\_unit = 0.04} bits per lookup epoch, interpreted in the identity-vs-mixture posterior-mass/PML sense rather than as a free-floating scalar \\
Replay hook & \texttt{plan\_id} + digest-bound \texttt{plan\_spec\_id} identify the accountant and the evidence needed to recompute the 20-step total from the declared retry / fallback schedule \\
State scope & \texttt{state\_decl\_id} is either absent because the plan proves the conditional step bound directly, or present to say which cache / routing / score state surfaces were included in that proof \\
Adjunct checkpoints & Any separation-failure guard appears explicitly as $(0,\rho)$ rather than being hidden inside the main $\eps$ sum, so later comparisons can tell whether the claim changed additively or multiplicatively \\
Owner routing & Anonymity~A owns the theorem, Synthesis~21 owns the conditional-step certificate packet, Synthesis~18 owns drift semantics for replay/state ids, and the Operational / State papers own the concrete mechanism evidence \\
\bottomrule
\end{tabularx}
\caption{A minimal public horizon-claim card. This is the smallest public answer later terse papers can quote directly when the live issue is whether a horizon-$Q$ receipt really declared an honest accountant, replay hook, and state boundary.}
\label{tab:minimal-horizon-card}
\end{table}

This card is intentionally non-decorative: it pins one threat window, one regime choice, one per-unit budget, one replay hook discipline, and one optional state / checkpoint boundary, without silently re-deriving the theorem or the underlying mechanism packet.

\paragraph{One tiny composition drill.}
If a receipt declares Regime~II with $Q=20$ uses and a per-use budget of $0.04$ bits, then the main horizon total is
\[
20\cdot 0.04 = 0.80\ \text{bits}.
\]
If the same receipt also carries one explicit separation-failure adjunct $\rho = 2^{-12}$, the honest public claim is therefore ``$0.80$ bits with an additional $(0,2^{-12})$ checkpoint'' rather than a bare ``20-use'' slogan.
A later terse paper may quote exactly that packet and move on, unless the live issue is how the conditional-step proof was replayed or which state surfaces were included.

\paragraph{A compact horizon-accounting first-reopen-owner card.}
If a later terse paper inherits one previously cited horizon-accounting shortcut and that shortcut is no longer exact, the first honest reopen is the first already-owned note that controls the changed field family rather than the widest neighboring review note. The card below keeps that fail-closed answer public and compact.

\begin{table}[t]
\centering
\small
\begin{tabularx}{0.98\linewidth}{@{}p{0.33\linewidth}p{0.22\linewidth}Y@{}}
\toprule
\textbf{If the stale horizon shortcut drifted in \dots} & \textbf{First honest reopen owner} & \textbf{Why stop there first}\\
\midrule
Threat-window unit, horizon boundary, or reset semantics & Synthesis~4 & Horizon accounting stays honest only after the repetition unit itself is fixed again.\\
Receipt tuple fields, replay-hook presence, or the packet shape that carries the accountant & Synthesis~12 & Before asking whether the accountant is right, later notes must know they are still talking about the same receipt-facing object.\\
Digest-bound \texttt{plan\_spec\_id} / \texttt{state\_decl\_id} sameness or the drift class saying whether state had to be declared & Synthesis~18 & This is the exact owner of replay/state equivalence and of the refresh-versus-replay-versus-recertify cut.\\
Regime choice, horizon total, or whether one explicit $(0,\rho)$ checkpoint is being stated additively rather than hidden in prose & Synthesis~19 & This note owns the receipt-facing accounting packet itself once the imported hooks are fixed.\\
Conditional-step certificate shape, conditioning interface, or per-step evidence packet for Regime~II & Synthesis~21 & Adaptive horizon claims stay honest only if the conditional-step certificate packet is still the same.\\
Concrete retry schedule, fallback evidence, or the rerun script proving the stated horizon total & Operational~A/B & Once the accountant packet is fixed, the next narrower live question is the operational replay evidence rather than horizon prose.\\
Concrete cache / routing / score state that must be included for the repeated-use claim to remain true & State~1 & Stateful repeats belong to the state witness tier, not to the generic accounting rulebook.\\
The theorem root showing why the one-vs-mixture horizon inequality composes at all & Anonymity~A & If the packet and evidence are fixed and only the proof root is live, reopen the theorem rather than the surrounding review stack.\\
\bottomrule
\end{tabularx}
\caption{A compact stale-shortcut routing card for horizon accounting. It names the first already-owned owner that regains authority once a previously cited horizon packet stops being exact.}
\label{tab:first-reopen-horizon}
\end{table}

\paragraph{One tiny first-reopen-owner drill.}
Suppose a successor keeps the same \texttt{tw\_id}, the same Regime~II declaration, the same digest-bound \texttt{plan\_spec\_id}, and the same \texttt{state\_decl\_id}, but changes the published per-step scalar from $0.04$ bits to $0.05$ bits or folds the explicit $\rho=2^{-12}$ checkpoint into untyped prose. The first honest reopen is still Synthesis~19, because the live issue is the horizon-accounting packet itself once its imports are fixed. If those scalars stay fixed but the conditioning interface or evidence object changes, the first honest reopen moves rightward to Synthesis~21. If both the packet and certificate stay fixed but the repetition unit changes from one lookup epoch to one hourly batch, the first honest reopen moves left to Synthesis~4. And if the same horizon slogan now depends on cache or routing churn that was previously outside scope, the first honest reopen moves to Synthesis~18 together with State~1, because the repeated-use claim has become stateful rather than purely accountant-local.\cite{series:threatwindows,series:receiptitems,series:planstateequiv,series:condcert,series:operationalA,series:operationalB,series:state1,series:oddsinflation}

\paragraph{Exports.}
We export (i) a receipt-facing definition of a \emph{horizon claim}, (ii) a small set of conservative editorial rules (when to state ``$Q\times$'' and when not to), and (iii) a checklist connecting horizon claims to replay plans (Operational series) and state declarations (State series).

\section{Horizon claims are objects, not prose}
Receipts in this archive are intentionally \emph{source-only}: they do not ship full tooling and proofs.
Therefore, when a receipt claims a horizon-$n$ budget, it must expose enough structure for an auditor to replay the accountant.
The key principle is: \emph{publish what you account, and account for what you publish.}

\begin{definition}[Horizon claim (receipt-facing)]\label{def:horizonclaim}
Fix a threat-window declaration (TW) naming the repetition unit (per lookup, per epoch, per day) and reset semantics.\cite{series:threatwindows}
A \emph{horizon claim} for a receipt line item consists of:
\begin{enumerate}[leftmargin=*]
\item a \textbf{unit} and \textbf{window binding} (TW-ID),
\item a \textbf{per-unit budget} $b$ (e.g., per-lookup MaxL / posterior-mass/PML cap),
\item an \textbf{accountant} $\mathsf{A}$ that maps $(b,n)$ to a cumulative budget $b_{\mathrm{tot}}$, and
\item a \textbf{replay hook} committing to how $\mathsf{A}$ is recomputed (\texttt{plan\_id} and digest-bound \texttt{plan\_spec\_id}).\cite{series:receiptitems,series:planstateequiv}
\end{enumerate}
The published scalar \texttt{budget\_value} is then interpreted as $b_{\mathrm{tot}}$ for the declared unit and window.
\end{definition}



\begin{remark}[How a horizon claim appears in a line item]
The receipt tuple schema (Synthesis~12) already has the needed hooks.\cite{series:receiptitems}
A horizon claim is realized by:
(i) binding the unit and reset semantics via \texttt{tw\_id},
(ii) declaring the usage assumption (often \texttt{usage}=$Q$ or a rate-limit clause), and
(iii) committing the \emph{accountant} and its evidence through the digest-bound \texttt{replay\_hook} (\texttt{plan\_id}, \texttt{plan\_spec\_id}).\cite{series:planstateequiv}
For Regime~II, the plan spec should in particular identify the conditioning interface and per-step evidence objects (Synthesis~21), so an auditor can replay the conditional-step inequality rather than assume independence.\cite{series:condcert}
\end{remark}

\begin{remark}[Why the replay hook matters]
Without a digest-bound plan spec, two releases can both say ``horizon-$n$'' but silently use different accounting rules (e.g., naive $n\eps$ vs checkpoint composition; or different retry policies), making public comparisons meaningless.
Operational~A/B and the worked example show what the replay plan looks like in practice.\cite{series:operationalA,series:operationalB,series:workedexample}
\end{remark}

\paragraph{Units.}
Throughout, we measure $\eps$ in \emph{bits} so that $2^{\eps}$ is the multiplicative identity-vs-mixture posterior-mass/PML cap.
(This matches the series convention used by Anonymity~A and the endpoint bridge.)\cite{series:oddsinflation}

\begin{proposition}[Adaptive sequential composition for one-vs-mixture budgets]\label{prop:adaptivecomp}
Let $K$ be the hidden key/identity and let $X_{1:n}$ be an interactive sequence of published witnesses (possibly chosen adaptively from past witnesses and state).
Write $P_k$ for the law given $K=k$ and $P=\sum_{j}\pi(j)P_j$ for the mixture under a fixed prior $\pi$.
Assume that for each step $t$ and each transcript prefix $x_{<t}$ with $P(x_{<t})>0$,
\[
P_k(X_t\in E\mid x_{<t})\ \le\ 2^{\eps_t}\,P(X_t\in E\mid x_{<t})\ +\ \delta_t
\qquad\text{for all measurable }E\text{ and all }k.
\]
Then the joint transcript satisfies the composed one-vs-mixture bound
\[
P_k(A)\ \le\ 2^{\sum_{t=1}^n\eps_t}\,P(A)\ +\ \sum_{t=1}^n \delta_t\,2^{\sum_{j>t}\eps_j}
\qquad\text{for all events }A.
\]
The additive slack is suffix-amplified because a failure at step $t$ is still carried through the multiplicative budget of later steps.  Heterogeneous receipt accountants must therefore preserve transcript order, not merely the multiset of per-step rows.
In particular, when all $\delta_t=0$ the realized posterior-mass inflation / PML (identity-vs-mixture log factor) composes additively: $\ell(X_{1:n})\le\sum_t\eps_t$.
\end{proposition}

\begin{remark}[Ownership of proofs]
The proof is the same hockey-stick / posterior-mass sequential composition argument used in Anonymity~A.\cite{series:oddsinflation}
This note records the \emph{receipt-facing condition}: you may claim a horizon-$n$ budget without independence if you commit (via the replay hook) to a plan that certifies the per-step conditional bound above.
\end{remark}

\begin{remark}[Separation-failure adjuncts compose as $\delta$]
If a primary-path surface is guarded by a separation manifest (Synthesis~22) and publishes a separation-failure upper bound $\rho$ (Synthesis~23), then the primary claim contributes a checkpoint $(0,\rho)$.
Receipt-side horizon accounting treats this exactly like any other $\delta$ term: it does not add to the $\varepsilon$ sum and it adds at most $\rho$ to the final additive slack (or appears as one step with $\delta_t=\rho$ in Proposition~\ref{prop:adaptivecomp}).\cite{series:primarysepmanifest,series:probsep}
In particular, a horizon-$n$ claim that composes per-step budgets $\{(\varepsilon_t,\delta_t)\}$ and carries one separation guard becomes
$(\sum_t\varepsilon_t,\ \rho\,2^{\sum_t\varepsilon_t} + \sum_t \delta_t\,2^{\sum_{j>t}\varepsilon_j})$ if the separation failure is a pre-transcript checkpoint, or
$(\sum_t\varepsilon_t,\ \rho + \sum_t \delta_t\,2^{\sum_{j>t}\varepsilon_j})$ if the separation guard is checked only at the final release boundary.
The receipt must name which timing it uses; it must not hide that choice inside a generic ``add $\rho$'' sentence.
\end{remark}

\section{Three composition regimes and what receipts may say}
This rulebook separates three common regimes.
A receipt should \emph{name} which regime it is in; otherwise it should publish only per-unit budgets.

\subsection{Regime I: independent repeats (rare, but clean)}
If the repeated witnesses are conditionally independent given the secret (or satisfy the independence conditions declared by the tiered-observation interface), then additive composition is often justified.\cite{series:contractobjects,series:tieredobs}
This is the cleanest case, but it is fragile for DHTs because routing tables, caches, and adversarial adaptivity create shared state.

\subsection{Regime II: adaptive sequential composition (default for retries)}
If the mechanism is used interactively, with later choices depending on earlier prefixes, then independence should not be assumed.
The receipt-grade condition is a \emph{per-step conditional} one-vs-mixture bound; under that condition, horizon composition follows by adaptive sequential composition (\cref{prop:adaptivecomp}). (certificate interface: Synthesis~21).\cite{series:condcert}
Operational~A/B show how to force retries and fallback decisions into auditable schedules so the per-step claim is meaningful, and Anonymity~A owns the underlying composition proof.\cite{series:operationalA,series:operationalB,series:oddsinflation}

\subsection{Regime III: stateful repeats (caches, scoring, routing churn)}
For anonymous DHTs, the secret can influence persistent state (routing-table updates, cache hits, score changes), and that state can later influence which witnesses are exposed.
In this regime, a naive ``$Q\times$'' claim is usually false unless the state surface is explicitly budgeted.
The State series treats state evolution as a first-class witness tier and exports the declarations needed for honest accounting.\cite{series:state1}

\begin{remark}[Editorial rule: no horizon-$Q$ without an explicit regime]
A receipt may publish a horizon-$Q$ (or ``$Q\times$'') total only if it declares which regime it is using:
(i) an explicit conditional-independence declaration justifying product-form repeats (Regime~I),
(ii) a per-step conditional budget that can be replayed (Regime~II; \cref{prop:adaptivecomp}), or
(iii) a digest-bound state declaration (\texttt{state\_decl\_id}) that enlarges the witness family so the per-step conditional bound is actually true in the presence of caches, routing churn, scoring, or other shared state (Regime~III).\cite{series:planstateequiv,series:state1}
Otherwise, the receipt should publish per-unit budgets only and leave horizon claims to a later, explicitly certified accountant.
\end{remark}

\section{Research warranted: dependence-aware equalization without bloat}
The series currently takes a conservative editorial stance: if state is not controlled and declared, do not claim $Q\times$ composition.
Two research directions can reduce conservatism while staying receipt-grade:
\begin{enumerate}[leftmargin=*]
\item \textbf{Contraction-based accounting.}
If secret influence contracts over time (e.g., via random-walk delegation, cover mixing, or a spectral gap), then horizon growth can be a discounted sum rather than linear.
The backbone spectral paper provides the overlay-specific distinguishability identities used by several budget papers.\cite{series:spectral}
A receipt-grade interface would expose the contraction factor as a knob with a dual certificate (mirroring the TV dual-certificate pattern).\cite{series:knapsacktransport}
\item \textbf{Stop-time normal forms for adaptive deadlines.}
Retries and deadlines often induce stopping-time effects.
The stop-time padding addendum provides normal forms (deadline tails, mixture separations) that can be used to keep accounting honest under optional stopping.\cite{series:stoptimeaddendum}
A compact ``optional-stopping accountant'' could then be referenced by both runtime receipts (Operational series) and evaluation receipts (Evaluation~1).
\end{enumerate}
These are intentionally stated as interface targets: once the owning mechanism paper exports a knob+certificate for contraction or stop-time effects, this note becomes the canonical place to say how it composes.

\section{Checklist for authors and auditors}
A paper (or receipt) that wants to state a horizon budget should satisfy:
\begin{enumerate}[leftmargin=*]
\item \textbf{TW declared:} unit, window, reset semantics (Synthesis~4).
\item \textbf{Per-unit budget exported:} posterior-mass/PML / MaxL / hockey-stick (mechanism paper + endpoint bridge).
\item \textbf{Accountant named:} which theorem/recipe is being used (Anonymity~A; State series; Boss Fight budgets).
\item \textbf{Replay hook bound:} \texttt{plan\_id} + \texttt{plan\_spec\_id} for auditors (Synthesis~12; Operational~B; worked example).
\item \textbf{State scope honest:} either independence is justified, or \texttt{state\_decl\_id} binds which state surfaces are included (State series; Synthesis~18).
\end{enumerate}

\begin{thebibliography}{99}\small

\bibitem{series:contractobjects}
Anonymous.
\newblock Anonymity Contract Objects for Anonymous DHTs: A Practical Interface for Proof-Carrying Claims.
\newblock Synthesis~0, The-One-True-Archive (2026).

\bibitem{series:receiptitems}
Anonymous.
\newblock Receipt Line Items for Anonymity Claims: A Minimal Tuple Schema for Published Budgets.
\newblock Synthesis~12, The-One-True-Archive (2026).

\bibitem{series:threatwindows}
Anonymous.
\newblock Threat Windows for Anonymous DHT Deployments.
\newblock Synthesis~4, The-One-True-Archive (2026).

\bibitem{series:oddsinflation}
Anonymous.
\newblock Posterior-Mass Inflation Anonymity: Pointwise Maximal Leakage, Composition, and the True-Odds Boundary.
\newblock Anonymity series Paper~A, The-One-True-Archive (2026).

\bibitem{series:operationalA}
Anonymous.
\newblock Schedule-Only Retries for Anonymous DHT Lookups.
\newblock Operational series Paper~A, The-One-True-Archive (2026).

\bibitem{series:operationalB}
Anonymous.
\newblock Auditable Trace Privacy: Replay Plans, PTL Evidence, and Receipt-Grade Monitoring.
\newblock Operational series Paper~B, The-One-True-Archive (2026).

\bibitem{series:workedexample}
Anonymous.
\newblock Worked Example: Receipt Interlock and Replay of a Tiered AnonDHT Claim.
\newblock Synthesis~17, The-One-True-Archive (2026).

\bibitem{series:tieredobs}
Anonymous.
\newblock Tiered Observation Vectors for Anonymous-DHT Receipts.
\newblock Synthesis~15, The-One-True-Archive (2026).

\bibitem{series:planstateequiv}
Anonymous.
\newblock Plan and State Equivalence: Digest-Bound Replay Hooks for Receipts.
\newblock Synthesis~18, The-One-True-Archive (2026).

\bibitem{series:state1}
Anonymous.
\newblock State-Dependent Anonymity for Anonymous DHT Deployments.
\newblock AnonDHT State series Paper~1, The-One-True-Archive (2026).

\bibitem{series:condcert}
Anonymous.
\newblock Conditional Per-Step Budget Certificates: A Receipt-Grade Interface for Adaptive Horizon Claims.
\newblock Synthesis~21, The-One-True-Archive (2026).

\bibitem{series:spectral}
Anonymous.
\newblock Spectral Anonymity and Optimal Distinguishing Bounds for Random-Walk Delegation in (Possibly Directed) Overlays.
\newblock Published backbone paper (2026-01-23), The-One-True-Archive.

\bibitem{series:stoptimeaddendum}
Anonymous.
\newblock Stop-Time Padding Addendum, Max-Mixture Separations and Tail-Sign Deadline Normal Forms.
\newblock Published backbone paper (2026-01-26), The-One-True-Archive.

\bibitem{series:knapsacktransport}
Anonymous.
\newblock Optimal Poset-Feasible Padding: Knapsack, Transport, and Dual Certificates for TV Privacy.
\newblock Published backbone paper (2026-01-29), The-One-True-Archive.



\bibitem{series:primarysepmanifest}
Anonymous.
\newblock Primary-Path Separation Manifests: Guard Objects for Non-Collusion Assumptions in Reader-Privacy Receipts.
\newblock Synthesis~22, The-One-True-Archive (2026).

\bibitem{series:probsep}
Anonymous.
\newblock Probabilistic Separation Budgets: Separation Failure as a Receipt-Grade $(0,\rho)$ Adjunct.
\newblock Synthesis~23, The-One-True-Archive (2026).

\end{thebibliography}

\end{document}
